\section{Derivación de la Parametrización Reducida de Grover}
\label{ap:grover_reducido}

La idea central es que, bajo las hipótesis habituales del algoritmo,
la evolución no necesita seguir de manera independiente las \(N\) amplitudes del vector de
estado. Debido a la simetría entre estados marcados y no marcados, basta con describir la
dinámica mediante dos parámetros: una amplitud común para los estados marcados y otra
para los no marcados.

Sea \(W \subseteq \{0,\dots,N-1\}\) el conjunto de estados marcados por el oráculo, con
\(M = |W|\). En la versión estándar del algoritmo, el estado inicial es la superposición
uniforme
\[
\ket{s}=\frac{1}{\sqrt{N}}\sum_{x=0}^{N-1}\ket{x},
\]
y el oráculo actúa cambiando el signo de las amplitudes asociadas a los estados de \(W\).

Bajo estas condiciones, todos los estados marcados tienen inicialmente la misma amplitud,
y lo mismo ocurre con todos los estados no marcados. Además, tanto el oráculo como el
operador de difusión preservan esa simetría: el primero afecta por igual a todos los estados
marcados, y el segundo transforma cada amplitud únicamente en función del promedio global.
Por ello, durante toda la evolución ideal del algoritmo, las amplitudes pueden agruparse en
solo dos clases:
\[
a_x^{(t)}=
\begin{cases}
u_t, & x\in W,\\[4pt]
v_t, & x\notin W,
\end{cases}
\]
donde \(u_t\) representa la amplitud común de los estados marcados en la iteración \(t\),
y \(v_t\) la amplitud común de los no marcados.

Esta misma idea puede expresarse en términos de una base reducida del subespacio relevante.
Definiendo los vectores normalizados
\[
\ket{w}=\frac{1}{\sqrt{M}}\sum_{x\in W}\ket{x},
\qquad
\ket{r}=\frac{1}{\sqrt{N-M}}\sum_{x\notin W}\ket{x},
\]
el estado inicial uniforme se descompone como
\[
\ket{s}
=
\frac{1}{\sqrt{N}}\sum_{x=0}^{N-1}\ket{x}
=
\sqrt{\frac{M}{N}}\,\ket{w}
+
\sqrt{\frac{N-M}{N}}\,\ket{r}.
\]
En consecuencia, la evolución ideal de Grover queda confinada al subespacio generado por
\(\{\ket{w},\ket{r}\}\), y cualquier estado del algoritmo puede escribirse como
\[
\ket{\psi_t}=\alpha_t \ket{w}+\beta_t \ket{r}.
\]

La relación entre ambas representaciones es directa. Como \(\ket{w}\) distribuye
uniformemente la amplitud sobre los \(M\) estados marcados y \(\ket{r}\) sobre los
\(N-M\) restantes, se tiene
\[
u_t=\frac{\alpha_t}{\sqrt{M}},
\qquad
v_t=\frac{\beta_t}{\sqrt{N-M}}.
\]
Por tanto, seguir la evolución en términos de \((\alpha_t,\beta_t)\) o en términos de
\((u_t,v_t)\) es equivalente.

\subsection*{Acción del Oráculo}

El oráculo de Grover cambia el signo de las amplitudes correspondientes a los estados
marcados y deja intactas las demás. Por tanto, si antes del oráculo el estado está descrito por
\((u_t,v_t)\), después del oráculo las amplitudes quedan como
\[
u_t'=-u_t,
\qquad
v_t'=v_t.
\]

\subsection*{Acción del Operador de Difusión}

El operador de difusión viene dado por
\[
D = 2\ket{s}\!\bra{s} - I,
\]
y puede interpretarse como una inversión de cada amplitud respecto al promedio global.
Si \(\mu_t\) denota el promedio de amplitud después de aplicar el oráculo, entonces la difusión
transforma cada componente según la regla
\[
x \mapsto 2\mu_t - x.
\]

Después del oráculo, el promedio de amplitud es
\[
\mu_t
=
\frac{1}{N}
\left(
\sum_{x\in W}(-u_t)+\sum_{x\notin W}v_t
\right)
=
\frac{-Mu_t+(N-M)v_t}{N}.
\]

Aplicando ahora la difusión:

Para un estado marcado, cuya amplitud después del oráculo es \(-u_t\),
\[
u_{t+1}=2\mu_t-(-u_t)=2\mu_t+u_t.
\]
Sustituyendo \(\mu_t\),
\[
u_{t+1}
=
2\left(\frac{-Mu_t+(N-M)v_t}{N}\right)+u_t
=
\frac{N-2M}{N}u_t+\frac{2(N-M)}{N}v_t.
\]

Para un estado no marcado, cuya amplitud después del oráculo es \(v_t\),
\[
v_{t+1}=2\mu_t-v_t.
\]
Sustituyendo nuevamente \(\mu_t\),
\[
v_{t+1}
=
2\left(\frac{-Mu_t+(N-M)v_t}{N}\right)-v_t
=
-\frac{2M}{N}u_t+\frac{N-2M}{N}v_t.
\]

En consecuencia, una iteración completa de Grover queda descrita por la recurrencia
\[
\begin{bmatrix}
u_{t+1}\\
v_{t+1}
\end{bmatrix}
=
\begin{bmatrix}
\frac{N-2M}{N} & \frac{2(N-M)}{N}\\[6pt]
-\frac{2M}{N} & \frac{N-2M}{N}
\end{bmatrix}
\begin{bmatrix}
u_t\\
v_t
\end{bmatrix}.
\]

\subsection*{Condiciones Iniciales y Reconstrucción del Estado}

Como el algoritmo parte de la superposición uniforme, en \(t=0\) todas las amplitudes son
iguales a \(1/\sqrt{N}\). Por tanto,
\[
u_0=v_0=\frac{1}{\sqrt{N}}.
\]

A partir de la recurrencia anterior, cada iteración puede resolverse actualizando únicamente
los parámetros \(u_t\) y \(v_t\), sin necesidad de aplicar el oráculo y la difusión sobre las
\(N\) amplitudes del vector de estado. Si en algún momento se requiere reconstruir el estado
completo, basta asignar
\[
a_x^{(t)}=
\begin{cases}
u_t, & x\in W,\\[4pt]
v_t, & x\notin W.
\end{cases}
\]

\subsection*{Alcance de la Reducción}

La parametrización reducida depende de una propiedad muy específica de Grover: la
evolución preserva la igualdad de amplitud dentro de las clases ``marcada'' y ``no marcada''.
Esa propiedad se cumple cuando el algoritmo parte de la superposición uniforme y el oráculo
solo introduce un cambio de fase sobre los estados marcados. Bajo estas hipótesis, la dinámica
efectiva queda contenida en un subespacio de dimensión dos.